\documentclass[10pt,a4paper]{amsart}
\usepackage[margin=19mm]{geometry}
\usepackage{amsmath,amssymb,mathtools}
\usepackage[scr=rsfs,cal=euler]{mathalfa}
\usepackage{xfrac}
\usepackage[unicode,hidelinks,bookmarks=false]{hyperref}
\newcommand{\ie}{i.e.\ }
\newcommand{\cf}{cf.\ }
\newcommand{\resp}{respectively\ }
\newcommand{\Subsection}{\subsection}
\providecommand{\nobreakdash}{\nobreak\hbox{-}}
\setlength{\emergencystretch}{2em}
\multlinegap0pt
\usepackage{bm}
\usepackage{bbm}
\usepackage{dsfont}
\usepackage{xspace}
\usepackage{cancel}
\usepackage{mathtools}
\usepackage{amsthm}
\makeatletter
\@ifundefined{theorem}{\newtheorem{theorem}{Theorem}}{}
\@ifundefined{lemma}{\newtheorem{lemma}[theorem]{Lemma}}{}
\makeatother
\newcommand{\A}{\mathbb{A}}
\title[$K$-Theory of Adelic and Rational Group $C^*$-algebras via G]{$K$-Theory of Adelic and Rational Group $C^*$-algebras via Generalized Winding Numbers}
\author{Wenqing Wu, Hang Wang}
\date{Source: arXiv:2509.00390v2 (CC BY 4.0)}
\begin{document}
\maketitle

\noindent\emph{Source and licence.} $K$-Theory of Adelic and Rational Group $C^*$-algebras via Generalized Winding Numbers. Version arXiv:2509.00390v2 (CC BY 4.0), distributed under the Creative Commons Attribution 4.0 International licence, as recorded in the arXiv licence metadata for this exact version and shown on the abstract page https://arxiv.org/abs/2509.00390v2. Full rights notice: licenses/arxiv-2509.00390-CC-BY-4.0.txt.\par\smallskip

\noindent\textit{Editorial scope.} Counted unit: the Lemma whose source label is \texttt{adelelift} in the article (source0.tex lines 683--685, source label \texttt{adelelift}), delivered together with the article's own complete original proof of it (source0.tex lines 686--779, one \texttt{proof} environment of 94 lines). No mathematics is written, repaired or re-derived: statement and proof are the article's own text line for line. Required context retained uncounted: lift (lines 421--427); lift8 (lines 449--451); gwd (lines 494--500); china (lines 562--574); gonggongwind (lines 625--634). Every definition, display and label of those lines is retained unchanged. Omitted from this package and not redistributed: 1--420, 428--448, 452--493, 501--561, 575--624, 635--682, 780--984. Nothing inside a retained excerpt is omitted or altered: every retained body is delivered exactly as the article's lines are written, so no omission receipt of any kind accompanies this package. Citations: the retained excerpts carry no citation command at all, so no bibliography entry of the article is transcribed and no reference is redistributed. Reference adaptation: none was needed and none was made; every reference inside the delivered unit resolves inside it, so no unresolved reference or citation is produced. Printed slips kept literal: no printed word, symbol, hypothesis or number of the counted statement or of its proof was altered. Layout and class: the article's own class is \texttt{article} with packages including geometry, amsthm, xcolor, amsmath, mathtools, amssymb, appendix, graphicx, placeins, abstract, enumerate, hyperref, biblatex; the wrapper is the lane's pinned \texttt{amsart} editorial wrapper at 10pt on A4, which restates the theorem-like environments the retained text needs and adds the article's own macro definitions that the retained text needs (\texttt{\textbackslash A}), with extra wrapper packages (bm, bbm, dsfont, xspace, cancel, mathtools). The delivered numbering is the unit's own. Assets: no retained span contains a figure environment, a graphics-inclusion command, a PDF-page inclusion, a table or a TikZ picture; no image file is copied into this package, the package declares no asset dependency and no image file is redistributed with it. Licence: version arXiv:2509.00390v2 of this article is distributed under the Creative Commons Attribution 4.0 International licence, as recorded in the arXiv licence metadata for this exact version and shown on the abstract page https://arxiv.org/abs/2509.00390v2. A full rights notice is delivered at licenses/arxiv-2509.00390-CC-BY-4.0.txt. Characters: every retained excerpt is pure ASCII, so the delivered engine, XeLaTeX with its default Latin Modern text font, reports no missing character.
\begin{lemma}\label{lift}
    Let the map \( p: \mathbb{R} \to S^1 \) be defined by \( p(x) = e^{2\pi ix} \). For any pre-periodic function \( f \in C(\mathbb{R}, S^1) \), and for any \( e \in p^{-1}(f(r)) \), there exists a function \( \tilde{f} \in C(\mathbb{R}, \mathbb{R}) \) such that the following conditions hold:
\begin{enumerate}
    \item \( \tilde{f}(r) = e \),
    \item \( p \circ \tilde{f} = f \).
\end{enumerate}
\end{lemma}

\begin{lemma}\label{lift8}
Let \( \tilde{f} \) be a lift of the pre-periodic function \( f \). Then, for \( \varepsilon \in [0, \frac{1}{8}) \), if the pre-periodic function \( g \) satisfies \( |g - f| < \varepsilon \), then there exists a unique lift \( \tilde{g} \) of \( g \) satisfies \( |\tilde{g} - \tilde{f}| < \frac{\varepsilon}{4}\).
\end{lemma}

\begin{theorem}\label{gwd}
f is a pre-periodic function, then the following limit exists and is a rational number:
$$
\lim_{x\to+\infty}\frac{\tilde{f}(x)-\tilde{f}(-x)}{2x}.
$$
This limit is called the generalized winding number of \(f\), denoted by \(\#f\).
\end{theorem}

\begin{lemma}\label{china}
    Let \( A_f \) denote \(\prod\limits_{p \text{ prime}} \mathbb{Q}_p
 \), which is the restricted product of all \( \mathbb{Q}_p \),  \( e_p(N) \) denote the exponent of the prime \( p \) in a positive integer N's prime factorization, i.e. 
\(
N = \prod\limits_{p \text{ prime}} p^{e_p(N)},
\)
and \(K_N:=\prod\limits_{p \text{ prime}}p^{e_p(N)}\mathbb{Z}_p\). 
Then there exists the following decomposition:

$$
\A_f=\bigsqcup_{\alpha \in \mathbb{Q}/N\mathbb{Z}} \alpha+K_N.
$$
\end{lemma}

\begin{theorem}\label{gonggongwind}
    Let \( g \in C(\mathbb{A}/\mathbb{Q}, S^1) \), \( x_f \in \mathbb{A}_f \), then the associated element in \( C(\mathbb{R}, S^1) \) given by  
\[
g_{x_f}(x_\infty) := g(\{x_\infty, x_f\})
\]  
is a pre-periodic function. Futhermore, for any \( x_f \in \mathbb{A}_f\), we have \(\#g_{x_f}=\#g_{0_f}\), where \(0_f\) is the embedding of 0 in \(\A_f\). Therefore, we can define the generalized winding number of \( g \in C(\mathbb{A}/\mathbb{Q}, S^1) \) by
$$
\#g:=\#g_{0_f}.
$$
\end{theorem}

\begin{lemma}\label{adelelift}
     For \( g \in C(\A/\mathbb{Q}, S^1) \) with \( \#g = 0 \), there exists \( \tilde{g} \in C(\A/\mathbb{Q}, \mathbb{R}) \) such that \( g = p \circ \tilde{g} \), where \(p(x)=e^{2\pi ix}\in C(\mathbb{R}).\)
\end{lemma}

\begin{proof}
    Let \( g_{0_f}(x_\infty) = g(\{x_{\infty}, 0_f\}) \) , which is a pre-periodic function on \( \mathbb{R} \). Following Lemma \ref{lift}, there exists \( \widetilde{g_{0_f}} \in C(\mathbb{R}, \mathbb{R}) \) such that \( p(\widetilde{g_{0_f}}) = g_{0_f} \).

Due to uniform continuity of $g$, we can find a $\delta > 0$ such that 
\[
|x - y|_{\A} < \delta \quad \Rightarrow \quad |g(x) - g(y)| < \frac{1}{8}.
\]
 Next, we select a suitable positive integer \( N \) such that the diameter of the set $K_N$ is less than $\delta$, which implies that
\[
\lvert g_{x_f}(x_\infty) - g_{0_f}(x_\infty)\rvert= \lvert g(\{x_{\infty}, x_f\}) - g(\{x_{\infty}, 0_f\})\rvert < \frac{1}{8},\ \forall x_f\in K_N.
\]

According to Lemma \ref{lift8}, there exists a unique lift \(\widetilde{g_{x_f}}\) of \( g_{x_f} \), satisfying:
\[
|\widetilde{g_{x_f}}(x_\infty) - \widetilde{g_{0_f}}(x_\infty)| < \frac{1}{32}.
\]

As stated in Lemma \ref{china},
\[
\A_f = \bigsqcup_{\alpha\in \mathbb{Q}/N\mathbb{Z}} \left( \alpha + K_N\right),
\]
then it follows that for any \( x_f \in \A_f \), there exists \( \alpha \in \mathbb{Q}/N\mathbb{Z} \) such that \( x_f + \alpha \in K_N \). 

For \( g \in C(\mathbb{A}/\mathbb{Q},S^1) \), we have:
\[
g_{x_f}(x_\infty) = g_{x_f + \alpha}(x_\infty+\alpha).
\]

Thus, by Lemma \ref{lift}, the function \( g_{x_f} \) has a unique lift at \( 0_{\infty} \) with value \( \widetilde{g_{x_f + \alpha}}(0_\infty+\alpha) \), which is denoted as \( \widetilde{g_{x_f}} \).

Define the map \( \widetilde{g}: \{x_{\infty}, x_f\} \mapsto \widetilde{g_{x_f}}(x_\infty) \) from \( \A \) to \( \mathbb{R} \). It is evident that $g=e^{2\pi i \tilde{g}}.$ Now we need to show that $\tilde{g}$ is continuous and $\mathbb{Q}$-periodic. To prove $\tilde{g}$ is continuous, we need to consider its continuity on $\mathbb{R}$ and $\A_f$. Since each $\widetilde{g_{x_f}}$ is the lift of a pre-periodic function (and hence a uniformly continuous function), $\tilde{g}$ is uniformly continuous with respect to the variable $x_\infty$. By the proof of Theorem \ref{gonggongwind}, $g$ is uniformly continuous, then for any $0<\varepsilon<\frac{1}{8}$, we can choose a $\delta >0$ such that:
\begin{align}
    1&, \lvert x_f-y_f\rvert_{\A_f}<\delta \Rightarrow x_f-y_f\in K_N, \forall x_f,\ y_f\in \mathbb{A}_f.\\
    2&,\lvert x_f-y_f\rvert_{\A_f}<\delta \Rightarrow \sup\lvert g_{x_f}-g_{y_f}\rvert<\varepsilon.
\end{align}
For the $x_f,y_f$ above, we have
\[
\lvert \widetilde{g_{x_f}}(x_\infty)-\widetilde{g_{y_f}}(x_\infty)\rvert<\frac{\varepsilon}{4},
\]
according to Lemma \ref{lift8}. Therefore, $\tilde{g}$ is uniformly continuous for both $x_\infty$ and $x_f$.

Finally, we prove that \( \widetilde{g} \in C(\A / \mathbb{Q}, \mathbb{R}) \), i.e. for any \( x_f \in \A_f, \alpha\in\mathbb{Q} \), it holds that:
\[
\widetilde{g}(\{x_\infty, x_f\}) = \widetilde{g}(\{x_\infty + \alpha, x_f + \alpha\}).
\]

With the positive integer $N$ given above, we divide the rational number $\alpha$ into the sum of $\alpha_1$ and $\alpha_2$, where \( \alpha_1 \in N\mathbb{Z} \) and \( \alpha_2 \in \mathbb{Q} /N\mathbb{Z} \). Let us further assume that $x_f\in \beta +K_N$ for some $\beta\in \mathbb{Q}/N\mathbb{Z}$.

Part 1: For the integer $N$ given above, we have
\[
|g_{x_f}(x_\infty) - g_{x_f}(x_\infty + kN)| < \frac{1}{8}, \ \forall x_\infty \in \mathbb{R}, k \in \mathbb{Z}.
\]
In view of Theorem \ref{gwd}, it is possible to find an integer $n_1$ such that:
\[
|\widetilde{g_{x_f}}(x_\infty) - \widetilde{g_{x_f}}(x_\infty + kN)-k\cdot n_1| \leq \frac{1}{32}.
\]
However, the generalized winding number of \(g_{x_f}\) is zero, which equals to $\frac{n_1}{N}.$ It follows that 
\[
|\widetilde{g_{x_f}}(x_\infty) - \widetilde{g_{x_f}}(x_\infty + kN)| \leq \frac{1}{32}.
\]
Additionally, 
\[
g_{x_f + kN}(x_\infty) = g(\{x_\infty, x_f + kN\}) = g(\{ x_\infty - kN, x_f \})=g_{x_f}(x_{\infty}-kN).
\]
So \( \widetilde{g_{x_f+kN}}(x_\infty ) \) and \( \widetilde{g_{x_f}}(x_\infty-kN) \)  differ by an integer.

Note that \(x_f\) and\( x_f+kN\) both belong to $\beta +K_N$, so \[|\widetilde{g_{x_{f}}}(x_{\infty})-\widetilde{g_{x_{f}+kN}}(x_{\infty})|<\frac{1}{32},\]while 
\[
|\widetilde{g_{x_{f}}}(x_{\infty})-\widetilde{g_{x_{f}}}(x_{\infty}+kN)|<\frac{1}{32}.
\]
Hence, for every \(\ k\in \mathbb{Z},x_{\infty}\in \mathbb{R},x_{f}\in \beta+K_N\)
\begin{align*}
    |\widetilde{g_{x_{f}+kN}}(x_{\infty})&-\widetilde{{g}_{x_{f}}}(x_{\infty}-kN)|<\frac{1}{16}.
\end{align*}
Thus,
\begin{align*}
    \widetilde{g_{x_{f}+kN}}&(x_{\infty})=\widetilde{{g}_{x_{f}}}(x_{\infty}-kN).
\end{align*}
We can replace \( kN \) with \( \alpha_1 \) to complete the first part of the proof.

Part 2: From the definition of \( \tilde{g} \),  for  \( x_f \in\beta+ K_N \), we have\[\widetilde{g_{x_f}}(0_{\infty})=\widetilde{g_{x_f-\beta}}(-\beta).\]
However, \(\widetilde{g_{x_f}}(x_{\infty})-\widetilde{g_{x_f-\beta}}(x_{\infty}-\beta)\) is a lift of the function
\[g_{x_f}(x_{\infty})\cdot g_{x_f-\beta}(x_{\infty}-\beta)^{-1},\]
which is equal to 1 for all $x_{\infty}\in \mathbb{R}$. Thus,
\[\tilde{g}(\{x_\infty, x_f\})=\widetilde{g_{x_f}}(x_{\infty})=\widetilde{g_{x_f-\beta}}(x_{\infty}-\beta)=\tilde{g}(\{x_\infty-\beta, x_f-\beta\}).\]
Similarly, for $x_f-\beta\in K_N, \beta+\alpha_2\in \mathbb{Q}$, it can be proved that:
\[
\widetilde{g_{x_f-\beta}}(x_{\infty}-\beta)=\widetilde{g_{x_f+\alpha_2}}(x_{\infty}+\alpha_2).
\]
We arrive at the final conclusion:
\[
\tilde{g}(\{x_\infty, x_f\})=\tilde{g}(\{x_\infty+\alpha, x_f+\alpha\}), \forall \{x_\infty,x_f\}\in\A, \alpha\in \mathbb{Q}.
\]
\end{proof}

\end{document}
